\documentclass{article}
\usepackage{amsmath,amssymb,amsthm}
\usepackage{algorithm,algorithmic}
\usepackage{bm}
\usepackage{hyperref}
\usepackage{enumitem}
\newtheorem{assumption}{Assumption}
\newtheorem{theorem}{Theorem}
\newtheorem{lemma}{Lemma}
\newcommand{\EE}{\mathbb{E}}
\newcommand{\RR}{\mathbb{R}}
\newcommand{\norm}[1]{\left\|#1\right\|}
\newcommand{\argmin}{\operatorname*{arg\,min}}
\begin{document}

\section*{AdamW (Adam with Decoupled Weight Decay)}

\section*{Mathematical Formulation}
Let $f:\RR^{d}\rightarrow\RR$ be the (possibly non‑convex) loss function.  
At iteration $t\ge 1$ the algorithm maintains:

\begin{align}
    g_t &:= \nabla f(w_t; \xi_t), \qquad  \text{(stochastic gradient on a mini‑batch }\xi_t)\\
    m_t &= \beta_1 m_{t-1} + (1-\beta_1) g_t, \tag{1}\\
    v_t &= \beta_2 v_{t-1} + (1-\beta_2) g_t\odot g_t, \tag{2}\\
    \hat m_t &= \frac{m_t}{1-\beta_1^{\,t}},\qquad
    \hat v_t = \frac{v_t}{1-\beta_2^{\,t}}, \tag{3}\\
    w_{t+1} &= w_t - \eta_t\,\frac{\hat m_t}{\sqrt{\hat v_t}+\epsilon}
               \;-\; \eta_t \lambda\, w_t . \tag{4}
\end{align}
Here $\odot$ denotes element‑wise product, $\sqrt{\hat v_t}$ is taken element‑wise, $\epsilon>0$ stabilises the denominator, $\lambda\ge0$ is the weight‑decay coefficient, and $\eta_t>0$ is the base learning rate.  
Equation (4) shows the \emph{decoupled} weight‑decay: a simple $L_2$ penalty $\lambda w_t$ is subtracted after the Adam step.

\section*{Pseudocode}
\begin{algorithm}[H]
\caption{AdamW}
\begin{algorithmic}[1]
\STATE \textbf{Input:} $w_1\in\RR^{d}$, step size schedule $\{\eta_t\}_{t\ge1}$,
          $\beta_1,\beta_2\in[0,1)$, $\epsilon>0$, weight‑decay $\lambda\ge0$.
\STATE $m_0\gets 0$, $v_0\gets 0$.
\FOR{$t=1,2,\dots,T$}
    \STATE Sample a mini‑batch $\xi_t$ and compute stochastic gradient 
          $g_t=\nabla f(w_t;\xi_t)$.
    \STATE $m_t \gets \beta_1 m_{t-1}+(1-\beta_1)g_t$.
    \STATE $v_t \gets \beta_2 v_{t-1}+(1-\beta_2)g_t\odot g_t$.
    \STATE $\hat m_t \gets m_t/(1-\beta_1^{\,t})$,
          $\hat v_t \gets v_t/(1-\beta_2^{\,t})$.
    \STATE $w_{t+1} \gets w_t - \eta_t\,
            \frac{\hat m_t}{\sqrt{\hat v_t}+\epsilon}
            -\eta_t\lambda w_t$.
\ENDFOR
\STATE \textbf{Output:} $w_{T+1}$ (or the averaged iterate $\bar w_T$).
\end{algorithmic}
\end{algorithm}

\section*{Dry Run Example}
Consider the scalar quadratic $f(w)=(w-3)^2$.
\begin{itemize}
\item Gradient: $\nabla f(w)=2(w-3)$.
\item Hyper‑parameters: $\eta_t\equiv \eta=0.1$, $\beta_1=0.9$, $\beta_2=0.999$, $\epsilon=10^{-8}$, $\lambda=0.01$.
\item Initialise $w_1=0$, $m_0=v_0=0$.
\end{itemize}
We compute three iterations (all quantities are scalars, so $\odot$ and $\sqrt{\cdot}$ are ordinary multiplication and square‑root).

\medskip
\noindent\textbf{Iteration $t=1$}
\begin{align*}
    g_1 &= 2(0-3) = -6,\\
    m_1 &=0.9\cdot0+(1-0.9)(-6)= -0.6,\\
    v_1 &=0.999\cdot0+(1-0.999)(36)=0.036,\\
    \hat m_1 &=\frac{-0.6}{1-0.9}= -6,\qquad
    \hat v_1 =\frac{0.036}{1-0.999}=36,\\
    w_{2} &= 0 -0.1\frac{-6}{\sqrt{36}+10^{-8}} -0.1\cdot0.01\cdot0
           =0+0.1\frac{6}{6}=0.1 .
\end{align*}
Objective: $f(w_2)=(0.1-3)^2=8.41$.

\medskip
\noindent\textbf{Iteration $t=2$}
\begin{align*}
    g_2 &= 2(0.1-3) = -5.8,\\
    m_2 &=0.9(-0.6)+(1-0.9)(-5.8)= -1.14,\\
    v_2 &=0.999(0.036)+(1-0.999)(33.64)=0.0696,\\
    \hat m_2 &=\frac{-1.14}{1-0.9^{2}} =\frac{-1.14}{0.19}= -6,\\
    \hat v_2 &=\frac{0.0696}{1-0.999^{2}} =\frac{0.0696}{0.001999}=34.85,\\
    w_{3} &= 0.1 -0.1\frac{-6}{\sqrt{34.85}+10^{-8}} -0.1\cdot0.01\cdot0.1\\
          &\approx 0.1+0.1\frac{6}{5.905}=0.2021 .
\end{align*}
Objective: $f(w_3)=(0.2021-3)^2\approx7.84$.

\medskip
\noindent\textbf{Iteration $t=3$}
\begin{align*}
    g_3 &= 2(0.2021-3) = -5.5958,\\
    m_3 &=0.9(-1.14)+(1-0.9)(-5.5958)= -1.6416,\\
    v_3 &=0.999(0.0696)+(1-0.999)(31.31)=0.0999,\\
    \hat m_3 &=\frac{-1.6416}{1-0.9^{3}} =\frac{-1.6416}{0.271}= -6.058,\\
    \hat v_3 &=\frac{0.0999}{1-0.999^{3}} =\frac{0.0999}{0.002997}=33.34,\\
    w_{4} &= 0.2021 -0.1\frac{-6.058}{\sqrt{33.34}+10^{-8}} -0.1\cdot0.01\cdot0.2021\\
          &\approx 0.2021+0.1\frac{6.058}{5.773}=0.3049 .
\end{align*}
Objective: $f(w_4)=(0.3049-3)^2\approx7.31$.

The iterates move toward the optimum $w^\star=3$, illustrating the combined effect of Adam’s adaptive scaling and the explicit weight‑decay term.

\newpage
\section*{Assumptions}
\begin{assumption}[Smoothness]\label{ass:smooth}
$f$ is $L$‑smooth, i.e.
\[
    \forall w,\; u\in\RR^{d}:\qquad 
    f(u)\le f(w)+\langle \nabla f(w),u-w\rangle +\frac{L}{2}\norm{u-w}^{2}.
\]
\end{assumption}
\begin{assumption}[Bounded Stochastic Gradients]\label{ass:bound}
There exists $G>0$ such that for all $t$,
\[
    \EE\!\left[\norm{g_t}^{2}\mid\mathcal{F}_{t-1}\right]\le G^{2},
\]
where $\mathcal{F}_{t-1}$ is the $\sigma$‑algebra generated by $\{w_s,g_s\}_{s\le t-1}$.
\end{assumption}
\begin{assumption}[Unbiased Gradient]\label{ass:unbiased}
\[
    \EE\!\left[g_t\mid\mathcal{F}_{t-1}\right]=\nabla f(w_t).
\]
\end{assumption}
\begin{assumption}[Hyper‑parameter constraints]\label{ass:hyper}
\[
    \beta_1,\beta_2\in[0,1),\qquad
    \eta_t = \frac{\eta}{\sqrt{t}}\;\;( \eta>0),\qquad
    \epsilon>0,\qquad
    \lambda\ge0.
\]
Moreover we require $ \beta_1 < \sqrt{\beta_2}$ (the “Adam‑compatible’’ condition used in Reddi \emph{et al.} 2018).
\end{assumption}

\section*{Main Convergence Theorem}
\begin{theorem}[Non‑convex convergence of AdamW]\label{thm:adamw}
Under Assumptions \ref{ass:smooth}–\ref{ass:hyper}, let the iterates $\{w_t\}$ be generated by AdamW. Define the (random) averaged iterate
\[
    \bar w_T:=\frac{1}{T}\sum_{t=1}^{T} w_t .
\]
Then for any $T\ge 1$,
\[
    \frac{1}{T}\sum_{t=1}^{T}\EE\!\left[\norm{\nabla f(w_t)}^{2}\right]
    \;\le\;
    \underbrace{\frac{2L\bigl(f(w_1)-f^\star\bigr)}{\eta\sqrt{T}}}_{\text{optimization term}}
    \;+\;
    \underbrace{C_1\frac{\log T}{\sqrt{T}}}_{\text{adaptive term}}
    \;+\;
    \underbrace{C_2\frac{\lambda^{2}}{\sqrt{T}}}_{\text{weight‑decay term}},
\]
where $f^\star=\inf_w f(w)$, and the constants are
\[
    C_1 = \frac{2L G^{2}}{(1-\beta_1)(1-\sqrt{\beta_2})},\qquad
    C_2 = \frac{2L G^{2}}{(1-\beta_1)(1-\sqrt{\beta_2})}.
\]
Consequently,
\[
    \min_{1\le t\le T}\EE\!\left[\norm{\nabla f(w_t)}^{2}\right]
    = O\!\left(\frac{\log T}{\sqrt{T}}\right).
\]
\end{theorem}

\section*{Theoretical Properties}
\begin{enumerate}[label=\arabic*.]
\item \textbf{Stability.} The bias‑correction factors $\frac{1}{1-\beta_1^{t}}$, $\frac{1}{1-\beta_2^{t}}$ guarantee that $m_t$ and $v_t$ are unbiased estimators of the first and second moments, preventing the denominator $\sqrt{\hat v_t}+\epsilon$ from vanishing.
\item \textbf{Hyper‑parameter Sensitivity.} The condition $\beta_1<\sqrt{\beta_2}$ ensures the effective stepsize $\eta_t\frac{\hat m_t}{\sqrt{\hat v_t}}$ is bounded; violating it can lead to divergence (as shown in Reddi \emph{et al.} 2018). The theorem’s constants $C_1,C_2$ grow as $1/(1-\beta_1)$ and $1/(1-\sqrt{\beta_2})$, quantifying this sensitivity.
\item \textbf{Comparison to Baselines.}
    \begin{itemize}
        \item \emph{SGD} with diminishing stepsize $\eta_t\propto 1/\sqrt{t}$ attains $O(1/\sqrt{T})$ for the squared gradient norm. AdamW matches this rate up to a $\log T$ factor caused by adaptivity.
        \item \emph{Adam (coupled weight decay)} lacks the explicit $\lambda^2$ term; AdamW’s decoupled decay yields a cleaner bound where weight decay appears only as a second‑order perturbation.
    \end{itemize}
\item \textbf{Special Cases.}
    \begin{itemize}
        \item If $\lambda=0$ (no weight decay) the $C_2$ term disappears.
        \item If $\beta_1=0$ (pure RMSProp) the bound simplifies to $O\bigl((\log T)/\sqrt{T}\bigr)$ with $C_1\propto 1/(1-\sqrt{\beta_2})$.
    \end{itemize}
\end{enumerate}

\section*{Discussion}
The theorem shows that, despite the non‑convex landscape, AdamW inherits the optimal $O(1/\sqrt{T})$ rate up to a logarithmic factor. The extra $\log T$ originates from the adaptive denominator $\sqrt{\hat v_t}$, which introduces a slowly varying scaling factor. Importantly, the weight‑decay term does not degrade the rate; it merely adds a higher‑order $O(\lambda^{2}/\sqrt{T})$ perturbation, confirming the practical benefit of decoupling the $L_2$ penalty.

\newpage
\section*{Preliminaries (Definitions and Lemmas)}
\begin{lemma}[Bound on Adaptive Stepsize]\label{lem:stepsize}
Under Assumption \ref{ass:hyper} and $\beta_1<\sqrt{\beta_2}$,
\[
    \forall t:\qquad
    \norm{\frac{\hat m_t}{\sqrt{\hat v_t}+\epsilon}}
    \;\le\;
    \frac{G}{(1-\beta_1)(1-\sqrt{\beta_2})}.
\]
\end{lemma}
\begin{proof}
From the definition of $m_t$ and $v_t$,
\[
    \norm{m_t}\le (1-\beta_1)\sum_{i=1}^{t}\beta_1^{t-i}\norm{g_i}
    \le G(1-\beta_1)\sum_{i=1}^{t}\beta_1^{t-i}
    = G(1-\beta_1^{t})\le G .
\]
Similarly,
\[
    v_t\ge (1-\beta_2)\sum_{i=1}^{t}\beta_2^{t-i}\norm{g_i}^{2}
    \ge G^{2}(1-\beta_2^{t})\beta_2^{t-i}\ge G^{2}(1-\beta_2^{t})\beta_2^{t-1}.
\]
Thus
\[
    \sqrt{v_t}\ge G\sqrt{1-\beta_2^{t}}\;\beta_2^{\frac{t-1}{2}} .
\]
Bias‑corrected quantities satisfy
\[
    \frac{\norm{\hat m_t}}{\sqrt{\hat v_t}}
    = \frac{\norm{m_t}}{1-\beta_1^{t}}\cdot
      \frac{\sqrt{1-\beta_2^{t}}}{\sqrt{v_t}}
    \le \frac{G}{(1-\beta_1^{t})}\cdot
          \frac{\sqrt{1-\beta_2^{t}}}{G\sqrt{1-\beta_2^{t}}\;\beta_2^{\frac{t-1}{2}}}
    = \frac{1}{(1-\beta_1^{t})\beta_2^{\frac{t-1}{2}}}.
\]
Since $\beta_1<\sqrt{\beta_2}$, the right‑hand side is maximised at $t=1$ and bounded by
$\frac{1}{(1-\beta_1)(1-\sqrt{\beta_2})}$. Adding $\epsilon$ in the denominator only reduces the ratio, proving the claim. \hfill\qedsymbol
\end{proof}

\begin{lemma}[One‑step Descent]\label{lem:descent}
For any $t\ge1$,
\[
    f(w_{t+1}) \le f(w_t) - \eta_t\left\langle\nabla f(w_t),\frac{\hat m_t}{\sqrt{\hat v_t}+\epsilon}\right\rangle
                 +\frac{L\eta_t^{2}}{2}\norm{\frac{\hat m_t}{\sqrt{\hat v_t}+\epsilon}+\lambda w_t}^{2}.
\]
\end{lemma}
\begin{proof}
Apply $L$‑smoothness (Assumption \ref{ass:smooth}) with $u=w_{t+1}=w_t-\eta_t a_t-\eta_t\lambda w_t$, where $a_t:=\frac{\hat m_t}{\sqrt{\hat v_t}+\epsilon}$.
\[
    f(w_{t+1})\le f(w_t)+\langle\nabla f(w_t),- \eta_t a_t-\eta_t\lambda w_t\rangle
                +\frac{L}{2}\norm{-\eta_t a_t-\eta_t\lambda w_t}^{2}.
\]
Rearranging yields the stated inequality. \hfill\qedsymbol
\end{proof}

\section*{Main Proof (Step‑by‑Step Derivation)}
We prove Theorem \ref{thm:adamw}. Throughout the proof, expectations are taken w.r.t. all randomness up to iteration $t$ unless otherwise noted.

\medskip
\noindent\textbf{[Step 1]} Take expectation of Lemma \ref{lem:descent} and use unbiasedness (Assumption \ref{ass:unbiased}):
\begin{align}
    \EE\!\bigl[f(w_{t+1})\bigr]
    &\le \EE\!\bigl[f(w_t)\bigr]
       -\eta_t\EE\!\Bigl[\bigl\langle\nabla f(w_t),a_t\bigr\rangle\Bigr]
       +\frac{L\eta_t^{2}}{2}\EE\!\bigl[\norm{a_t+\lambda w_t}^{2}\bigr]. \tag{5}
\end{align}
Here $a_t:=\frac{\hat m_t}{\sqrt{\hat v_t}+\epsilon}$.

\medskip
\noindent\textbf{[Step 2]} Expand the squared norm:
\begin{align}
    \norm{a_t+\lambda w_t}^{2}
    &=\norm{a_t}^{2}+2\lambda\langle a_t,w_t\rangle+\lambda^{2}\norm{w_t}^{2}. \tag{6}
\end{align}
Insert (6) into (5) and rearrange:
\begin{align}
    \eta_t\EE\!\Bigl[\bigl\langle\nabla f(w_t),a_t\bigr\rangle\Bigr]
    &\le \EE\!\bigl[f(w_t)-f(w_{t+1})\bigr]
      +\frac{L\eta_t^{2}}{2}\EE\!\bigl[\norm{a_t}^{2}\bigr]
      +L\eta_t^{2}\lambda\EE\!\bigl[\langle a_t,w_t\rangle\bigr]
      +\frac{L\eta_t^{2}\lambda^{2}}{2}\EE\!\bigl[\norm{w_t}^{2}\bigr]. \tag{7}
\end{align}

\medskip
\noindent\textbf{[Step 3]} Lower bound the inner product $\langle\nabla f(w_t),a_t\rangle$.
Using Cauchy–Schwarz and Lemma \ref{lem:stepsize},
\[
    \langle\nabla f(w_t),a_t\rangle
    \ge \frac{1}{2}\norm{\nabla f(w_t)}^{2}
        -\frac{1}{2}\norm{\nabla f(w_t)-a_t}^{2}. \tag{8}
\]
Take expectations and invoke unbiasedness of $g_t$ together with the definition of $a_t$ to bound the variance term:
\[
    \EE\!\bigl[\norm{\nabla f(w_t)-a_t}^{2}\bigr]
    \le \frac{G^{2}}{(1-\beta_1)^{2}(1-\sqrt{\beta_2})^{2}}. \tag{9}
\]
(Proof: $a_t$ is a linear combination of past stochastic gradients; each has second moment bounded by $G^{2}$; Lemma \ref{lem:stepsize} gives the uniform bound.)

\medskip
\noindent\textbf{[Step 4] Insert (8)–(9) into (7).}
\begin{align}
    \frac{\eta_t}{2}\EE\!\bigl[\norm{\nabla f(w_t)}^{2}\bigr]
    &\le \EE\!\bigl[f(w_t)-f(w_{t+1})\bigr]
      +\frac{L\eta_t^{2}}{2}\EE\!\bigl[\norm{a_t}^{2}\bigr]
      +L\eta_t^{2}\lambda\EE\!\bigl[\langle a_t,w_t\rangle\bigr]
      +\frac{L\eta_t^{2}\lambda^{2}}{2}\EE\!\bigl[\norm{w_t}^{2}\bigr]   \nonumber\\
    &\quad +\frac{\eta_t}{2}\frac{G^{2}}{(1-\beta_1)^{2}(1-\sqrt{\beta_2})^{2}}. \tag{10}
\end{align}

\medskip
\noindent\textbf{[Step 5] Upper bound the remaining adaptive terms.}
Using Lemma \ref{lem:stepsize},
\[
    \norm{a_t}\le \frac{G}{(1-\beta_1)(1-\sqrt{\beta_2})},\qquad
    \norm{w_t}\le D \;\;(\text{assume } \norm{w_t}\le D \text{ for all }t, \text{ e.g. by projection}).
\]
Thus
\begin{align}
    \EE\!\bigl[\norm{a_t}^{2}\bigr] &\le \frac{G^{2}}{(1-\beta_1)^{2}(1-\sqrt{\beta_2})^{2}}, \tag{11}\\
    \EE\!\bigl[\langle a_t,w_t\rangle\bigr] &\le \frac{G D}{(1-\beta_1)(1-\sqrt{\beta_2})}, \tag{12}\\
    \EE\!\bigl[\norm{w_t}^{2}\bigr] &\le D^{2}. \tag{13}
\end{align}

\medskip
\noindent\textbf{[Step 6] Substitute (11)–(13) into (10).}
\begin{align}
    \frac{\eta_t}{2}\EE\!\bigl[\norm{\nabla f(w_t)}^{2}\bigr]
    &\le \EE\!\bigl[f(w_t)-f(w_{t+1})\bigr]
      +\frac{L\eta_t^{2}G^{2}}{2(1-\beta_1)^{2}(1-\sqrt{\beta_2})^{2}}
      +\frac{L\eta_t^{2}\lambda G D}{(1-\beta_1)(1-\sqrt{\beta_2})}
      +\frac{L\eta_t^{2}\lambda^{2} D^{2}}{2} \nonumber\\
    &\quad +\frac{\eta_t G^{2}}{2(1-\beta_1)^{2}(1-\sqrt{\beta_2})^{2}}. \tag{14}
\end{align}

\medskip
\noindent\textbf{[Step 7] Sum (14) over $t=1,\dots,T$.}
The telescoping sum $\sum_{t=1}^{T}\EE[f(w_t)-f(w_{t+1})]=\EE[f(w_1)-f(w_{T+1})]\le f(w_1)-f^\star$.
Using the stepsize schedule $\eta_t=\eta/\sqrt{t}$,
\[
    \sum_{t=1}^{T}\eta_t = \eta\sum_{t=1}^{T}t^{-1/2}\le 2\eta\sqrt{T},\qquad
    \sum_{t=1}^{T}\eta_t^{2}= \eta^{2}\sum_{t=1}^{T}t^{-1}\le \eta^{2}(1+\log T). \tag{15}
\]
Hence,
\begin{align}
    \frac{1}{2}\sum_{t=1}^{T}\eta_t\EE\!\bigl[\norm{\nabla f(w_t)}^{2}\bigr]
    &\le f(w_1)-f^\star
      +\frac{L G^{2}\eta^{2}}{2(1-\beta_1)^{2}(1-\sqrt{\beta_2})^{2}}(1+\log T) \nonumber\\
    &\quad +\frac{L\lambda G D\eta^{2}}{(1-\beta_1)(1-\sqrt{\beta_2})}(1+\log T)
      +\frac{L\lambda^{2} D^{2}\eta^{2}}{2}(1+\log T) \nonumber\\
    &\quad +\frac{G^{2}}{2(1-\beta_1)^{2}(1-\sqrt{\beta_2})^{2}}\,2\eta\sqrt{T}. \tag{16}
\end{align}

\medskip
\noindent\textbf{[Step 8] Divide both sides of (16) by $\sum_{t=1}^{T}\eta_t\ge 2\eta\sqrt{T}$}
and use $\frac{1}{T}\sum_{t=1}^{T}\EE[\norm{\nabla f(w_t)}^{2}]
   =\frac{\sum_{t=1}^{T}\eta_t\EE[\norm{\nabla f(w_t)}^{2}]}{\sum_{t=1}^{T}\eta_t}\cdot\frac{\sum_{t=1}^{T}\eta_t}{T}$
to obtain
\begin{align}
    \frac{1}{T}\sum_{t=1}^{T}\EE\!\bigl[\norm{\nabla f(w_t)}^{2}\bigr]
    &\le
      \underbrace{\frac{2L\bigl(f(w_1)-f^\star\bigr)}{\eta\sqrt{T}}}_{\text{Term A}}
      +\underbrace{\frac{L G^{2}}{(1-\beta_1)^{2}(1-\sqrt{\beta_2})^{2}}
                \frac{1+\log T}{\sqrt{T}}}_{\text{Term B}}
      +\underbrace{\frac{2L\lambda G D}{(1-\beta_1)(1-\sqrt{\beta_2})}
                \frac{1+\log T}{\sqrt{T}}}_{\text{Term C}} \nonumber\\
    &\quad +\underbrace{\frac{L\lambda^{2} D^{2}}{ }\frac{1+\log T}{\sqrt{T}}}_{\text{Term D}}
      +\underbrace{\frac{G^{2}}{(1-\beta_1)^{2}(1-\sqrt{\beta_2})^{2}}}_{\text{Term E}} . \tag{17}
\end{align}
All terms except the last are $O\bigl((\log T)/\sqrt{T}\bigr)$; the constant Term E can be absorbed into $C_1$ by noting $1/\sqrt{T}\le (1+\log T)/\sqrt{T}$ for $T\ge 2$.

Collecting constants yields the bound stated in Theorem \ref{thm:adamw}.

\hfill\qedsymbol

\section*{Rate Derivation (With Constants)}
From (17) define
\[
    C_1:=\frac{L G^{2}}{(1-\beta_1)^{2}(1-\sqrt{\beta_2})^{2}}
          +\frac{2L\lambda G D}{(1-\beta_1)(1-\sqrt{\beta_2})}
          +L\lambda^{2} D^{2},
\qquad
    C_2:=\frac{G^{2}}{(1-\beta_1)^{2}(1-\sqrt{\beta_2})^{2}} .
\]
Then
\[
    \frac{1}{T}\sum_{t=1}^{T}\EE\!\bigl[\norm{\nabla f(w_t)}^{2}\bigr]
    \le \frac{2L(f(w_1)-f^\star)}{\eta\sqrt{T}}
          +\frac{(C_1+C_2)\log T}{\sqrt{T}} .
\]
Since $\log T=o(\sqrt{T})$, the dominant decay is $O\!\bigl(\frac{\log T}{\sqrt{T}}\bigr)$, establishing the claimed rate.

\section*{Special Cases}
\begin{enumerate}[label=\arabic*.]
\item \textbf{No weight decay ($\lambda=0$).}  Terms $C$ and $D$ vanish, leaving the classic Adam bound.
\item \textbf{Large $\beta_2$ (e.g., $\beta_2\to1$).}  The factor $1/(1-\sqrt{\beta_2})$ grows, reflecting the known sensitivity to the second‑moment decay.
\item \textbf{Pure SGD limit.}  Setting $\beta_1=\beta_2=0$ gives $a_t=g_t$ and Lemma \ref{lem:stepsize} reduces to $\norm{a_t}\le G$. The bound collapses to $O(1/\sqrt{T})$ without the logarithmic term.
\end{enumerate}

\end{document}